\subsection{Training Setup}

We use the following hyperparameters for SAC training. Learning rates are set to $3 \times 10^{-4}$ for policy, dual critics, and temperature networks. Discount factor is $\gamma = 0.99$. Batch size is 256. Target network exponential moving average coefficient is $\tau = 0.005$. Warm-start transitions (initial random exploration before RL updates) is 2000. Replay buffer capacity is 100k transitions. Target entropy for automatic temperature tuning is $H_{\text{target}} = -7$. Estimated training duration is approximately 20 hours on 4× RTX 6000 Ada GPUs with 8 parallel environments for sample efficiency.

The training loop operates as follows. An episode begins by sampling a water-body type from $\{\text{lake}, \text{river}, \text{coastal}, \text{open\_ocean}\}$ and correspondingly sampling mission parameters such as current magnitude and direction, wave characteristics, and wind. An initial position within valid spawn zones is sampled. Battery SoC is initialized to a random value in $[0.7, 1.0]$ to encourage learning across battery states. At each timestep, the state $s_t$ is constructed from vehicle kinematics, battery state, environmental measurements, inferred water-body type, frontier map, and sensor history. The policy samples a discrete mode $a_1$ and continuous actions $a_{2:7}$. Stochastic sensor activation probabilities are sampled to produce a binary active sensor vector. Hydrodynamic forces are computed from thrust, steering, and water current/wave models. Physics integration advances vehicle position and velocity. Energy consumption combines propulsive energy (function of thrust and velocity) plus sensor energy (function of active sensors and timestep duration). Battery SoC is decremented. Information gain is computed from active sensors using proxy metrics. Coverage map is updated. Reward is computed from the composite function. Transition $(s_t, a_t, r_t, s_{t+1}, \text{done})$ is stored. If replay buffer size exceeds warm-start threshold, SAC updates are applied to critics, policy, and temperature networks. Episodes terminate when battery SoC reaches zero (failure), collision occurs, coverage reaches 90\% (success), or maximum timesteps elapse.

Curriculum learning proceeds in four phases. Phase 1 (episodes 0--500) uses simple lake environments without currents, full battery starts, and clear weather. Phase 2 (episodes 500--1500) introduces mild currents ($0.3$--$0.5$ m/s), partial battery starts (70--100\% SoC), and occasional fog. Phase 3 (episodes 1500--3000) adds realistic waves (0.5--1.5 m), varied water types, and diverse weather. Phase 4 (episodes 3000+) enables full complexity including extreme conditions, challenging current patterns, and battery degradation effects. Curriculum learning accelerates convergence to stable policies by starting with simpler dynamics.

\subsection{Simulation Environments}

Primary simulation uses the Maritime Robotics Simulator built on Gazebo physics engine with UUV Simulator plugins providing hydrodynamic models. The simulator implements 6-DoF vehicle dynamics with nonlinear drag, current-induced forces, and wave effects via spectral wave models. Bathymetric maps represent different water bodies. Battery models track capacity degradation over charge cycles and temperature-dependent efficiency. Multi-sensor simulation includes Radar return statistics using K-distribution clutter models, Lidar range returns with wave-induced noise, optical camera image synthesis under fog and rain, and water quality sensor readings with spatial variation.

Four primary scenarios evaluate diverse conditions. Lake exploration uses a 5 km$^2$ calm water area with minimal current (0.1 m/s maximum), wave height less than 0.2 m, 90\% coverage target. River navigation uses a 3 km linear transect with steady 1.5 m/s flow, obstacles (islands, shallows), current-assisted drift opportunities. Coastal mapping uses a 10 km$^2$ area with significant waves (0.5--1.5 m), tidal currents (0.3--0.8 m/s), and varied water depth affecting sensor performance. Open ocean mission uses 20 km$^2$ with large waves (2--4 m), variable currents, long-endurance 12+ hour target emphasizing energy conservation.

\subsection{Baselines}

We compare against five baseline approaches. Frontier-based exploration greedily navigates to nearest unexplored frontier using standard RRT* path planning, all sensors always active, ignoring battery state. Lawn-mower coverage pattern follows fixed systematic sweeps with all sensors active, representative of standard maritime survey operations. Rule-based duty-cycle applies a hand-tuned schedule (sensors on 2 minutes, off 1 minute) and returns to base when battery SoC drops below 50\%, representative of operational practice. Energy-CPP uses SAC-based coverage path planning adapted from agricultural robotics, optimizes motion but uses fixed sensor scheduling. DDPG-ASV applies deterministic policy gradients to energy-constrained navigation, adapted from autonomous vehicle work, optimizes motion but ignores sensor duty-cycling.

\subsection{Evaluation Metrics}

Primary metrics directly address the mission objectives. Area mapped per Joule ($\text{m}^2/\text{J}$) measures coverage efficiency, the primary metric of interest. Mission success rate is the percentage of episodes achieving 90\% coverage with safe return (SoC $\geq 10\%$ at base). Average information gain sums sensor-weighted measurements accounting for quality degradation. Energy remaining at mission end (percentage SoC) indicates how much battery capacity was left unexploited.

Secondary metrics provide diagnostic insight. Coverage completion time measures hours to reach 90\% coverage target. Sensor utilization tracks percentage of mission duration each sensor remained active, revealing energy-saving strategies. Path efficiency is the ratio of area coverage distance to total distance traveled, detecting unnecessary detours. Safe return violations count episodes where the ASV returned with insufficient battery or failed to return at all.

Water-type specific metrics include per-environment coverage efficiency (lake, river, coastal, open ocean separately), adaptation quality measured as episode count to achieve 90\% of steady-state performance in new environment, and current exploitation effectiveness (for rivers, what percentage of drift opportunity was utilized for energy savings).

\subsection{Stress Test Scenarios}

Extreme current (3 m/s headwind) tests energy-conserving drift strategies. The policy should learn to minimize thrust and exploit any cross-current positioning to conserve energy, accepting slower coverage if necessary.

High wave state (4--6 m significant height) evaluates sensor reliability and collision avoidance. Lidar performance degrades in rough seas, camera performance is stable, Radar captures vessel echoes amidst sea clutter. The policy should learn sensor trade-offs.

Low battery start (30\% initial SoC) forces conservative exploration and early return trigger, testing safe return guarantees under tight constraints.

Sensor failures disable random sensors mid-mission (e.g., Lidar failure at episode 50\% completion). The policy should adapt coverage strategy, relying on remaining sensors.

Communication dropout (simulated loss of base station contact for extended periods) removes opportunity for route re-planning. The policy must make autonomous decisions to return despite uncertainty.

Multi-stressor scenarios combine effects such as high waves plus low battery plus sensor failure to evaluate policy robustness under realistic failure modes.

\subsection{Ablation Studies}

Water-body type conditioning ablates the water-body type one-hot vector (4D), removing explicit environment signal. Performance drop across environments reveals how much of the efficiency gain depends on water-type adaptation versus general learning.

Sensor duty-cycle ablation fixes sensor activation to always-on, testing whether the joint optimization benefit comes from motion-sensing coupling or just better motion planning. Expected result is moderate performance drop, confirming duty-cycle importance.

Discrete mode selection ablation uses flat continuous actions only (removing 4D discrete mode choice), testing whether interpretability constraint impacts performance. Expected result is slight performance drop with increased latency due to larger action space.

Reward component ablation removes individual reward terms (coverage-only, energy-only, safety-only, etc.) to quantify their contribution. Expected result is degraded performance without safety term, indicating its criticality.

These ablations validate that each component contributes meaningfully to the final performance.

\subsection{Deployment Specification}

Real-time inference on Jetson AGX Orin embedded system is critical for onboard ASV deployment. State vector construction (sensor fusion, metric computation) requires approximately 5ms on CPU. Policy forward pass (87D $\to$ 4D discrete $\to$ 7D continuous) requires approximately 3ms on embedded GPU. Action dispatch to motor controller and sensor manager requires approximately 2ms. Safety override computation (checking safe return threshold) requires approximately 1ms. Total latency budget is approximately 11ms, leaving 1ms margin for jitter within a 12ms control cycle (approximately 83 Hz capability, exceeding typical ASV 10 Hz control rates).

Inference on RTX 3080 desktop GPU achieves approximately 2ms policy forward pass, enabling 500 Hz simulation step rate for training and evaluation.

The lightweight policy network (250k parameters) fits within onboard memory constraints (less than 10 MB), enabling deployment on mid-range ASV platforms with limited compute (Jetson AGX Orin: 32 GB RAM, 8-core ARM CPU, 8× NV GPUs with 100 TFLOPs combined).
